% SPDX-License-Identifier: LPPL-1.3c
% Copyright (C) 2026 Christophe Jorssen
% Author and Current Maintainer: Christophe Jorssen <christophe.jorssen@gmail.com>
% LPPL maintenance status: maintained
% This file is part of luafemm. See LICENSE and LICENSES.md for its terms.
\documentclass[varwidth=160mm,border=3mm]{standalone}
\usepackage{pgfplots}
\pgfplotsset{compat=1.18}
\usetikzlibrary{femm}
\begin{document}
\tikzset{femm/materials/rotor iron/.style={mu r=200}}
\begin{tikzpicture}[
  femm/problem={xmin=-8,xmax=8,ymin=-8,ymax=8,mesh size=1,depth=10},
  femm/sliding band={name=time-band,inner radius=3,outer radius=4,
    segments=48,layers=1},scale=2]
  \filldraw[fill=gray!20,femm/region={material=rotor iron}]
    (0,0) ellipse[x radius=2.6,y radius=1.2];
  \draw[femm/region={material=copper,conductor=stranded,current density=5e7}]
    (-6,-3) rectangle (-5,3);
  \draw[femm/region={material=copper,conductor=stranded,current density=-5e7}]
    (5,-3) rectangle (6,3);
\end{tikzpicture}
\begin{tikzpicture}[
  femm/problem={sliding band=time-band,rotor angle=25},
  femm/state={name=rotor-time,sliding band=time-band,speed=0},scale=2]
  \pic {femm geometry}; \pic[femm/lines=13] {femm field};
  \node[above] at (0,8) {$t=0$};
\end{tikzpicture}
\par
\def\MotionResults{}
\foreach \Step in {1,...,20}{%
  \begin{tikzpicture}[
    femm/problem={previous state=rotor-time,time step=.001,
      motion=dynamic,inertia=1e-8,damping=1e-7},
    femm/state={name=rotor-time},scale=2]
    \femmmeasure[name=motion-\Step,quantity=force,method=band]
    \ifnum\Step=20
      \pic {femm geometry}; \pic[femm/lines=13] {femm field};
      \node[above] at (0,8) {$t=20\,\mathrm{ms}$};
    \fi
  \end{tikzpicture}%
  \xdef\MotionResults{\MotionResults,motion-\Step}%
}
\begin{tikzpicture}
\begin{axis}[width=8cm,height=5cm,grid=major,
  xlabel={Time (ms)},ylabel={Rotor angle (degrees)},
  x filter/.expression={1000*x}]
  \femmresultplot[results=\MotionResults,abscissa=time,component=angle,
    plot options={blue!65!black,thick,mark=*}]
\end{axis}
\end{tikzpicture}
\femmstate[name=rotor-saved,file=\jobname-rotor.lft]
\end{document}
